%%%% CAPÍTULO 1 - INTRODUÇÃO

\chapter{Introdução}\label{cap:introducao}

Mercados eletrônicos modernos são sistemas computacionais distribuídos nos quais a formação de preços depende de mensagens de oferta, cancelamento, execução e atualização de \gls{orderbook}, isto é, do livro de ofertas, que percorrem redes de baixa latência. Nesse cenário, a capacidade de receber dados de mercado, atualizar estado interno e produzir uma decisão em intervalos previsíveis tornou-se um requisito técnico central para aplicações de negociação algorítmica e, em especial, para \gls{hft}. A literatura de microestrutura descreve que a redução de latência altera a dinâmica competitiva dos participantes e influencia liquidez, descoberta de preços e comportamento do \gls{orderbook} \cite{HasbrouckSaar2013,Hendershott2011,Brogaard2014}.

Do ponto de vista de engenharia de computação, a dificuldade não está apenas em executar uma regra de negociação, mas em construir uma cadeia de processamento com baixo tempo de resposta, controle explícito de filas, representação adequada dos dados e integração entre software e hardware. CPUs de uso geral oferecem flexibilidade e ecossistema maduro, porém dependem de execução sequencial, hierarquia de memória, interrupções, escalonamento de sistema operacional e variabilidade temporal. \glspl{fpga}, por outro lado, permitem transformar partes do fluxo em circuitos dedicados, explorando paralelismo espacial e latência determinística no nível de ciclos de clock \cite{HarrisHarris2013,PattersonHennessy2017}.

Este trabalho investiga essa alternativa por meio de uma plataforma acadêmica de baixo custo baseada na Intel/Terasic DE10-Nano, que combina processador ARM no \gls{hps} com lógica programável Cyclone V. A contribuição principal está no uso de \gls{hls} com MATLAB HDL Coder para transformar uma estratégia de trading escrita em MATLAB em um bloco VHDL integrável a uma infraestrutura previamente preparada. Nessa abordagem, a comunicação ARM--FPGA, o fluxo de dados de mercado, o \gls{orderbook} simplificado, os testes e o projeto Quartus formam a base do ambiente; o ponto de variação passa a ser a estratégia, que pode ser redigida em MATLAB e regenerada para hardware.

\begin{figure}[htb]
\caption{Placa DE10-Nano utilizada no protótipo}
\label{fig:de10nano}
\centering
\includegraphics[width=0.72\textwidth]{DE10Nano}
\fonte{Terasic Technologies \cite{TerasicDE10Nano}.}
\end{figure}

\section{Considerações iniciais}\label{sec:consideracoesIniciais}

O termo \gls{hft} é empregado neste trabalho para representar sistemas automatizados que processam dados de mercado e emitem decisões em janelas temporais muito reduzidas. A implementação desenvolvida não se conecta a uma bolsa real, não executa ordens financeiras e não deve ser interpretada como sistema de produção. Ela é uma plataforma experimental para estudar os elementos computacionais que aparecem em uma cadeia de HFT: recepção de mensagens, normalização de eventos, estrutura de \gls{orderbook}, decisão algorítmica, medição de comportamento e integração com lógica programável.

A escolha por uma plataforma de baixo custo tem relevância acadêmica. Infraestruturas profissionais de HFT utilizam enlaces especializados, placas de rede de latência ultrabaixa, \textit{kernel bypass}, co-localização e FPGAs de alto desempenho. Esses recursos são financeiramente inacessíveis para a maioria dos laboratórios de graduação. A DE10-Nano, ilustrada na \autoref{fig:de10nano}, apesar de limitada em frequência, memória e recursos lógicos quando comparada a placas industriais, oferece um ambiente suficientemente rico para estudar os princípios de integração entre ARM, FPGA, barramentos, síntese de hardware e processamento determinístico \cite{TerasicDE10Nano,IntelCycloneV}. Esse caráter de baixo custo não é apenas qualitativo: a própria página de compra da Terasic lista a DE10-Nano em condição acadêmica por US\$190 \cite{TerasicDE10NanoPrice}. Para efeito de contraste, FPGAs de classe usada em trabalhos de HFT de alta performance podem custar ordens de grandeza a mais; por exemplo, o dispositivo AMD/Xilinx XCKU115-2FLVA1517E aparece listado pela Mouser por US\$10.290,57 \cite{MouserXCKU115Price}. Mesmo sendo uma comparação conservadora --- pois contrapõe uma placa completa a um único FPGA de alto desempenho --- ela já evidencia a diferença de orçamento entre uma plataforma didática e uma solução de classe profissional.

\section{Problema de pesquisa}\label{sec:problemaPesquisa}

O problema tratado neste TCC pode ser formulado da seguinte forma: como construir, em uma plataforma FPGA de baixo custo, uma infraestrutura que permita implementar estratégias de trading com menor esforço de descrição RTL, preservando determinismo, paralelismo em hardware e integração com um fluxo de dados de mercado? Essa pergunta combina duas dificuldades. A primeira é arquitetural: sistemas baseados apenas em CPU sofrem com processamento sequencial, overhead do sistema operacional, variabilidade de escalonamento e limitações do modelo von Neumann. A segunda é metodológica: escrever toda a lógica de decisão diretamente em VHDL aumenta a barreira de entrada para testar novas estratégias. A hipótese de trabalho é que separar infraestrutura estável e estratégia variável pode reduzir essa barreira sem abrir mão da execução em FPGA.

Do ponto de vista técnico, os requisitos derivados desse problema são: baixa variabilidade temporal no trecho em hardware; representação numérica fixa para preços e quantidades; processamento de \gls{orderbook} em estrutura de profundidade limitada; interface ARM--FPGA documentada; validação por simulação e execução em placa; e planejamento de métricas de latência, vazão, uso de recursos e produtividade na alteração de estratégias.

\section{Objetivos}\label{sec:objetivos}

Esta seção apresenta o objetivo geral e os objetivos específicos do trabalho.

\subsection{Objetivo geral}\label{subsec:objetivoGeral}

Desenvolver e documentar uma plataforma de HFT de baixo custo, baseada em ARM e FPGA, cujo objetivo principal é facilitar a criação de estratégias de trading em MATLAB e sua geração para hardware por meio do MATLAB HDL Coder, mantendo previamente preparados o fluxo de dados de mercado, o \gls{orderbook} simplificado, a comunicação ARM--FPGA e a infraestrutura de validação.

\subsection{Objetivos específicos}\label{subsec:objetivosEspecificos}

Os objetivos específicos são:

\begin{itemize}
  \item estruturar um ambiente no qual a estratégia de trading seja o principal bloco substituível pelo usuário;
  \item definir um modelo de estratégia em MATLAB compatível com HDL Coder, utilizando tipos sintetizáveis, entradas numéricas e saída de ação;
  \item automatizar o fluxo de geração VHDL da estratégia a partir do código MATLAB;
  \item definir o contrato de comunicação entre software ARM e lógica FPGA, incluindo registradores, ponteiros, semântica de filas e formato dos quadros;
  \item implementar e explicar o fluxo de geração e recepção de mensagens FAST sintéticas utilizado como alimentação de mercado;
  \item manter, no FPGA, um \gls{orderbook} por níveis de preço com profundidade fixa, ordenação por lado e cálculo de melhor compra, melhor venda, \textit{spread} e desbalanceamento;
  \item integrar a estratégia gerada por MATLAB HDL Coder ao motor VHDL por meio de uma interface estável;
  \item validar funcionalmente os componentes por testes C++, simulações VHDL e teste de fumaça entre gerador e receptor;
  \item propor métricas para avaliar latência, vazão, uso de recursos FPGA e esforço necessário para alterar uma estratégia.
\end{itemize}

\section{Justificativa}\label{sec:justificativa}

A relevância do trabalho decorre da interseção entre arquitetura de computadores, sistemas embarcados, redes e aplicações financeiras de baixa latência. Em HFT, pequenas variações de tempo de resposta podem modificar a ordem de chegada de mensagens e, consequentemente, a decisão disponível ao participante. A literatura mostra que atividades de negociação algorítmica e baixa latência afetam liquidez e descoberta de preços, ainda que os efeitos dependam do mercado, do período e da estratégia adotada \cite{Hendershott2011,Brogaard2014,HasbrouckSaar2013}. Portanto, estudar a infraestrutura computacional subjacente é relevante para compreender tanto ganhos de desempenho quanto riscos técnicos.

Do ponto de vista pedagógico, o trabalho força a integração de conteúdos normalmente vistos de forma separada. O receptor C++ exige programação de sistemas, sockets TCP, manipulação binária e interface com bibliotecas externas. O bloco FPGA exige VHDL, filas circulares, protocolos de barramento, controle de \textit{handshake} e testbenches. A etapa MATLAB HDL Coder introduz \gls{hls} e as restrições de escrever código algorítmico sintetizável.

A escolha por \gls{fpga} é justificada pelo tipo de carga de trabalho. Atualizações de \gls{orderbook}, empacotamento de quadros, cálculo de \textit{spread}, comparação de limiares e geração de ação são operações regulares, com largura fixa e baixa necessidade de memória dinâmica. Essas características são favoráveis à implementação em hardware dedicado. Em contrapartida, decodificação de protocolo, reconexão TCP, mensagens de erro e interação com sistema operacional permanecem no ARM, onde a flexibilidade de software é mais apropriada. Esse co-projeto evita uma solução puramente acadêmica desconectada das restrições práticas de depuração e integração.

Além da motivação técnica, há impacto educacional e econômico. Educacionalmente, o trabalho aproxima sistemas embarcados, arquitetura de computadores, mercado financeiro, HLS e validação experimental em um único projeto. Economicamente, a escolha por uma DE10-Nano e por ferramentas acadêmicas permite avaliar uma alternativa de baixo custo frente a soluções comerciais especializadas. Cientificamente, o projeto contribui com uma análise prática de como ferramentas modernas de HLS podem ser aplicadas a um domínio sensível a latência, mas sem exigir que cada nova estratégia seja escrita manualmente em RTL.

\section{Estrutura do trabalho}\label{sec:estruturaTrabalho}

O \autoref{cap:referencialTeorico} apresenta o referencial teórico utilizado para compreender HFT, \gls{orderbook}, FPGA, HLS, MATLAB HDL Coder, FAST, MMIO, Avalon-MM, filas circulares e métricas de desempenho. O \autoref{cap:trabalhos:relacionados} discute trabalhos relacionados em negociação algorítmica, aceleração em FPGA, HLS e processamento de \glspl{orderbook}. O capítulo de proposta descreve a plataforma desenvolvida, suas fronteiras, o fluxo de geração de estratégias e seus contratos de dados. O \autoref{cap:materialemetodos} detalha materiais, ferramentas, métodos de desenvolvimento e planejamento de métricas. O \autoref{cap:resultados} apresenta o sistema obtido, a modelagem, a implementação e os espaços reservados para resultados quantitativos. Por fim, o \autoref{cap:conclusoeseperspectivas} sintetiza as contribuições, limitações e perspectivas futuras.
